\section{Discussion}\label{sec:discussion}

\subsection{Where the delay was}

The delay removed from the transport was not where delay is usually sought. The
encoder took 1.2 to 1.5~ms and the loopback network nothing; two
frame periods were
held in the receiving process, one because its parser recognizes the
end of a frame
by the start of the next, and one because its output stage was
threaded by frame.
Neither is a property of H.264 or of RTP, and the packets on the
tether are the same
before and after. Stage-by-stage measurements in the literature find,
as here, that
encoding, the network and decoding need a few milliseconds
\cite{bachhuber2018,lv2026,georg2020}, while pipelines assembled from
standard parts
report 60 to 200~ms without assigning it to a stage
\cite{green2021,khalil2026,schimpe2022,testouri2025}; the 70~ms from
which we started
was one of these. We did not measure those pipelines and cannot say
that the same two
causes are at work in them. They are, however, properties of a widely
used decoder
when it is fed over RTP, and not of our sender, so that they are worth excluding
wherever such a delay is unexplained. For the ROV observed at 850~ms
with H.264 and
500~ms with MJPEG, the authors recommended cameras of low compression
\cite{kaknjo2018}, which needed more than ten times the bit rate
there; our result
suggests that the codec need not be given up for the sake of delay.

What remains, 3.8 to 8.3~ms, is the transport's share only. It
excludes the camera,
whose delay after exposure alone ranged from 2 to 82~ms in one survey
\cite{ubik2021}, and the display; the two were more than half of the
delay from scene
to display in an optimized chain \cite{bachhuber2018}. With the transport at a
quarter of a frame period, the camera and the display are now where the pilot's
delay is, and the capture path of Section~\ref{sec:capture}, which
discards all but
the newest buffer, addresses the part of it that software controls. Whether the
67~ms removed matters to a pilot is not settled by this work. Delays
of that size
lengthen fine manipulation \cite{mackenzie1993,nankaku2022}, the
shortest delay found
to degrade driving is 170~ms \cite{chen2007}, and for underwater
vehicles we found
only impressions \cite{bauer2025,doniec2010}. The saving is best
regarded as room in
the pilot's delay budget, which image processing can then spend.

\subsection{The visibility correction}

The correction belongs to the cheapest class of enhancement and uses its usual
operators; what it adds is that the parameters of those operators, and not the
pixels, follow their targets through a first-order filter. On a still scene
consecutive frames then differ by at most one level in fewer than
0.1\% of pixels,
and a step in brightness is followed within 1.5~s. That time constant
is a choice
between flicker and lag: a pilot who turns the lights on sees an over- or
under-corrected picture for a second or so. Filtering parameters
costs nothing per
pixel and cannot produce the ghosting of temporal filters on pixels
\cite{ancuti2012}, but it steadies only what is global, and the tile
mappings of the
local equalization are the least global part of it.

Its limits are those of its class. A global correction cannot undo a
loss that grows
with range \cite{berman2020}; in a scene that spans near and far, the
far part stays
blue and hazy. The red patch of the synthetic scene recovered only in
part, because
no gain restores a channel of which a quarter was transmitted without
also raising
its noise. On a flat, dark frame the standard deviation rose
3.9~times in spite of
the limit on gain, so that the correction amplifies sensor noise when there is
nothing else to amplify. And the evaluation is on synthetic frames
from a model that
is itself a simplification \cite{akkaynak2018}, a practice with known weaknesses
\cite{berman2020}; the constants of Table~\ref{tab:autovis} were set
on such frames
and a laptop camera and will change on footage from the vehicle.

On photographs from the water the correction came closer to the pictures that
observers preferred than the other methods of its cost did, and closer than our
implementations of two methods that took eighty and five hundred times as
long (Table~\ref{tab:uieb}). We read this with three reservations. The
references were chosen among the outputs of twelve methods, so that a method
scores well by resembling what those methods produce, and a correction that
stretches contrast and removes a cast moderately is of that kind; the grey-world
balance and the dark channel prior, which change colour most, lay furthest from
them. Fusion and the dark channel prior are our implementations: the second agrees
with the figure published for its authors' code, and the first, for which there
is none, was given the most favourable of the settings we tried. And photographs do not test what
the correction adds to its class, the filter on its parameters, since each is
the first frame of a video; flicker and lag remain evaluated on synthetic
sequences only. UCIQE ranked first two methods that lay
further from the references in signal-to-noise ratio than the unprocessed
photographs, in agreement with the criticisms of such measures (Section~\ref{sec:rw-enhancement}); we report them
in Table~\ref{tab:uieb} and do not argue from them. The correction lowered the
signal-to-noise ratio of a quarter of the photographs, which we have not yet
examined one by one. Since
unprocessed pictures serve object detectors best \cite{wang2024detectors}, the
correction is applied to the pilot's picture alone and the camera's frame is
published unaltered beside it.

\subsection{The frame bus}

The frame bus is the non-blocking write protocol in its form with
several buffers
\cite{kopetz1993}, mapped from a file and with a futex for waiting. It has that
protocol's central property, that the writer is not influenced by its
readers, and
therefore the property that matters on a vehicle: a program that reads frames,
however slow or faulty, cannot delay the pilot's picture. It does not have the
protocol's bound on the reader's repetitions, which rests on
guarantees of timing
that Linux does not give; a reader repeats a copy only when preempted for about
three frame periods, and then receives a newer frame. Against
transports that pass a
pointer \cite{pohnl2023} the bus pays a copy for each frame on each
side, 34.9~$\mu$s
to publish at 1280$\times$720 on a desktop processor, and in exchange needs no
daemon, no allocation in shared memory by readers, and no library: a
reader is a few
lines in any language that can map a file. For frames of this size at
30 frames per
second the copy is a small price; for many streams of much larger
frames it would
not be. The cost of the copy on the vehicle's processor has not been measured.

\subsection{Placement and the vehicle's processor}

The one measurement on a single-board computer shows where the
vehicle's processor
goes: of 0.91 of a core for the whole of electric-eye at
1280$\times$720, the H.264
encoder took 0.51, on a processor throttled by heat. Image
processing, capture and
the bus together are the smaller part. This is the reason for sending
a camera's own
H.264, which removes the encoder from the vehicle, and it has a price: without
decoded frames on the vehicle the operators run at the surface, and the pilot's
picture is encoded a second time there for the ground station. That
mode has been
tested with files and a loopback device and not yet with a camera that encodes.

\subsection{Limitations}

The software has not been run on a vehicle. The latency was measured
over loopback on
one laptop, without the tether's modems, the vehicle's camera or its
processor; the
processing costs on ARM are counts of instructions under emulation, which do not
represent memory behaviour, and one throttled run on a Raspberry
Pi~5; the camera was
a laptop's. The equivalence checks were run under emulation for both ARM
targets, and an emulator tests its own rendering of the instructions;
only the test
suites have run on an AArch64 processor. The visibility
correction has been evaluated on synthetic frames and on still photographs,
against two baselines of our own implementation, and not on video from the
water; no pilot has flown with it.
The next steps are accordingly the vehicle itself: the costs and the
latency on its
processor and across its tether, the delay of its camera and of the
display by an
optical method \cite{bachhuber2016}, the correction's constants on
footage from the
water, and a comparison of pilots' performance with and without the
correction and
the shortened transport.
